\chapter{\label{ch:SSH}Su-Schrieffer-Heeger model with multiorbital superconductivity}
As a first departure into the search for systems exhibiting topological 
superconductivity, we begin with a topologically non-trivial system.
A canonical topological system is the Su-Schrieffer-Heeger (SSH) model of a 
spinless, one-dimensional dimer lattice of fermions with staggered hopping. 
The model obeys a {\it chiral symmetry}, exhibiting two topological phases. 
In the topologically non-trivial phase, a pair of edge states emerges at 
the boundaries, characterised by a {\it winding number}. 
We will consider the effect of a spin-triplet, multiorbital attractive interaction within the unit cell, and a Coulomb repulsion to nearest-neighbours outside of the unit cell.

\section{\label{sec:SSH} Su-Schrieffer-Heeger model}
The Su-Schrieffer-Heeger model with additional spin degrees of freedom is given by the following Hamiltonian
\begin{multline}
    \hat{\mathcal H}_\text{SSH} =  
    \sum_{x=1}^{n_x} \sum_{\sigma\in\{\uparrow\downarrow\}} \sum_{m\in\{A,B\}} 
    \left[
-\mu \hat c_{x, \sigma, m}^\dagger \hat c_{x, \sigma, m}
\right]
+\sum_{\sigma\in\{\uparrow\downarrow\}} \left[V \hat c_{0, \sigma, m}^\dagger \hat c_{0, \sigma, m}\right]
    \\+
    \sum_{x=1}^{n_x} \sum_{\sigma\in\{\uparrow\downarrow\}}
    \left[
    -v \hat c_{x, \sigma, A}^\dagger \hat c_{x, \sigma, B}
    -w \hat c_{x, \sigma, B}^\dagger \hat c_{x+1, \sigma, A}
    +\text{h.c.}
    \right],
\label{eq:one-dimensional_SSH}
\end{multline}
where $\mu$ is the chemical potential, $v$ is the intracell, interorbital hopping parameter, $w$ is the intercell, interorbital hopping and $V$ is an impurity coupling strength. The model boundary conditions are periodic, and therefore, in the topological state, we expect edge modes to arise around the impurity.
In the trivial state, with $v>w$, the unit cell has the strucure displayed in Figure \ref{fig:SSH-unit-cell}.

We investigate the phases of the model in Figure \ref{fig:normal-SSH-phase-diagram}.

\createfigure{SSH-unit-cell}{SSH unit cell}

\createfigure{normal-SSH-phase-diagram}{Normal SSH phase diagram}

\createfigure{normal-SSH-LDOS-zero}{Normal SSH LDOS}

\createfigure{normal-SSH-LDOS-nonzero}{Normal SSH LDOS}

\section{\label{sec:spin-triplet_SSH} Non-unitary spin-triplet Su-Schrieffer-Heeger model}
The non-unitary spin-triplet extension of the Su-Schrieffer-Heeger model is given by
\begin{equation}
    \hat{\mathcal H}_\text{SSH-triplet} =  
    \hat{\mathcal H}_\text{SSH} +
    \sum_{x=1}^{n_x} \sum_{\sigma\in\{\uparrow,\downarrow\}}
    \left[
    -U_v \hat c_{x, \sigma, A}^\dagger \hat c_{x, \sigma, B}^\dagger \hat c_{x, \sigma, B} \hat c_{x, \sigma, A}
    +U_w \hat c_{x, \sigma, B}^\dagger \hat c_{x+1, \sigma, A}^\dagger \hat c_{x+1, \sigma, A} \hat c_{x, \sigma, B}
\right],
\label{eq:one-dimensional_SSH}
\end{equation}
where $U_v$ is the spin-triplet interaction and $U_w$ is a Coulomb repulsion.

The mean-field solution is
\begin{multline}
    \hat{\mathcal H}_\text{SSH-triplet}^\text{mf} =  
    \hat{\mathcal H}_\text{SSH} +
    \sum_{x=1}^{n_x} \sum_{\sigma\in\{\uparrow,\downarrow\}}
    \Bigg[
    \sum_{m\in\{A,B\}}\left[
    \phi_{x,\sigma, m} \hat c_{x, \sigma, m}^\dagger \hat c_{x, \sigma, m} \right]\\
    +\Phi_{x,\sigma}^v \hat c_{x, \sigma, A}^\dagger \hat c_{x, \sigma, B}
    +\Phi_{x,\sigma}^w \hat c_{x, \sigma, B}^\dagger \hat c_{x+1, \sigma, A}
    +\Delta_{x,\sigma}^v \hat c_{x, \sigma, A}^\dagger \hat c_{x, \sigma, B}^\dagger 
    +\Delta_{x,\sigma}^{v\dagger} \hat c_{x, \sigma, B} \hat c_{x, \sigma, A}\\
    +\Delta_{x,\sigma}^w \hat c_{x, \sigma, B}^\dagger \hat c_{x+1, \sigma, A}^\dagger 
    +\Delta_{x,\sigma}^{w\dagger}
    \hat c_{x+1, \sigma, A} \hat c_{x, \sigma, B}
\Bigg],
\label{eq:one-dimensional_SSH}
\end{multline}
where the Hartree, Fock and Gorkov fields are defined as
\begin{eqnarray}
\phi_{x,\sigma, m}&=& (U_v-U_w)\hat c_{x, \sigma, m}^\dagger \hat c_{x, \sigma, m}\\
\Phi_{x,\sigma}^v&=& -U_v \hat c_{x, \sigma, A}^\dagger \hat c_{x, \sigma, B}\\
\Phi_{x,\sigma}^w&=&\hat U_w c_{x, \sigma, B}^\dagger \hat c_{x+1, \sigma, A}\\
\Delta_{x,\sigma}^{v}&=& -U_v \hat c_{x, \sigma, B} \hat c_{x, \sigma, A}\\
\Delta_{x,\sigma}^{v\dagger}&=& U_w \hat c_{x, \sigma, A}^\dagger \hat c_{x, \sigma, B}^\dagger \\
\Delta_{x,\sigma}^{v}&=& -U_v \hat c_{x+1, \sigma, A} \hat c_{x, \sigma, B} \\
\Delta_{x,\sigma}^{w\dagger}&=& U_w \hat c_{x, \sigma, B}^\dagger \hat c_{x+1, \sigma, A}^\dagger 
\end{eqnarray}
